
\documentclass{article}
\usepackage{colortbl}
\usepackage{makecell}
\usepackage{multirow}
\usepackage{supertabular}

\begin{document}

\newcounter{utterance}

\centering \large Interaction Transcript for game `dond', experiment `coop', episode 3 with nemotron{-}3{-}ultra.
\vspace{24pt}

{ \footnotesize  \setcounter{utterance}{1}
\setlength{\tabcolsep}{0pt}
\begin{supertabular}{c@{$\;$}|p{.15\linewidth}@{}p{.15\linewidth}p{.15\linewidth}p{.15\linewidth}p{.15\linewidth}p{.15\linewidth}}
    \# & $\;$A & \multicolumn{4}{c}{Game Master} & $\;\:$B\\
    \hline

    \theutterance \stepcounter{utterance}  
    & & \multicolumn{4}{p{0.6\linewidth}}{
        \cellcolor[rgb]{0.9,0.9,0.9}{
            \makecell[{{p{\linewidth}}}]{
                \texttt{\tiny{[P1$\langle$GM]}}
                \texttt{Hrajete kooperativní vyjednávací hru, ve které se musíte dohodnout na tom, jak rozdělit sadu předmětů mezi sebou a druhým hráčem.} \\
\\ 
\texttt{Pravidla:} \\
\texttt{(a) Vy a druhý hráč obdržíte sadu předmětů. Každý z vás také obdrží tajnou funkcí hodnoty, která představuje, jak moc si ceníte jednotlivých typů objektů.} \\
\texttt{(b) Vyměňujete si zprávy s druhým hráčem, abyste se dohodli na tom, kdo dostane které předměty. Můžete poslat maximálně 5 zpráv každý nebo hru předčasně ukončit tím, že kdykoli podáte tajný návrh.} \\
\texttt{(c) Od každého z vás se vyžaduje, aby předložil tajný návrh označující položky, které chcete, formátované v hranatých závorkách takto: "{[}Návrh: <číslo> <název objektu>, <číslo> <název objektu>, <...>{]}"} \\
\texttt{(d) Pokud jsou vaše návrhy komplementární, tj. je dostatek předmětů na splnění obou návrhů, každý hráč obdrží skóre na základě součtu hodnot předmětů, které obdržel. Jinak oba dostanete nula bodů.} \\
\texttt{(e) Vaším cílem je maximalizovat součet vašeho skóre a skóre druhého hráče.} \\
\\ 
\texttt{Pojďme začít.} \\
\\ 
\texttt{Sada dostupných položek je:} \\
\\ 
\texttt{1 kleště, 2 cihly, 1 prsten, 1 hodinky} \\
\\ 
\texttt{Vaše tajná funkce hodnoty je:} \\
\\ 
\texttt{kleště: 1, cihla: 0, prsten: 2, hodinky: 7} \\
\\ 
\texttt{DŮLEŽITÉ: Vaše zprávy, pokud se nejedná o tajný návrh, jsou předány přímo druhému hráči, takže do nich nezahrnujte žádnou odpověď na pravidla ani text oznamující vaši zprávu. Chcete{-}li podat tajný návrh, použijte uvedený formát. Nepoužívejte hranaté závorky při komunikaci s druhým hráčem, jinak bude interpretován jako váš tajný návrh.} \\
            }
        }
    }
    & \\ \\

    \theutterance \stepcounter{utterance}  
    & \multicolumn{4}{p{0.6\linewidth}}{
        \cellcolor[rgb]{0.9,0.9,0.9}{
            \makecell[{{p{\linewidth}}}]{
                \texttt{\tiny{[P1$\rangle$GM]}}
                \texttt{Ahoj! Podívám se na naše předměty. Pro mě jsou nejdůležitější hodinky (7 bodů) a prsten (2 body). Kleště jsou 1 bod a cihly 0 bodů.} \\
\\ 
\texttt{Co je pro tebe důležité? Rád bychom to rozdělili tak, abychom obě získali maximum bodů.} \\
            }
        }
    }
    & & \\ \\

    \theutterance \stepcounter{utterance}  
    & & & \multicolumn{2}{p{0.3\linewidth}}{
        \cellcolor[rgb]{0.9,0.9,0.9}{
            \makecell[{{p{\linewidth}}}]{
                \texttt{\tiny{[GM$|$GM]}}
                \texttt{continue} \\
            }
        }
    }
    & & \\ \\

    \theutterance \stepcounter{utterance}  
    & & & \multicolumn{2}{p{0.3\linewidth}}{
        \cellcolor[rgb]{0.9,0.9,0.9}{
            \makecell[{{p{\linewidth}}}]{
                \texttt{\tiny{[GM$|$GM]}}
                \texttt{Ahoj! Podívám se na naše předměty. Pro mě jsou nejdůležitější hodinky (7 bodů) a prsten (2 body). Kleště jsou 1 bod a cihly 0 bodů.} \\
\\ 
\texttt{Co je pro tebe důležité? Rád bychom to rozdělili tak, abychom obě získali maximum bodů.} \\
            }
        }
    }
    & & \\ \\

    \theutterance \stepcounter{utterance}  
    & & \multicolumn{4}{p{0.6\linewidth}}{
        \cellcolor[rgb]{0.9,0.9,0.9}{
            \makecell[{{p{\linewidth}}}]{
                \texttt{\tiny{[GM$\rangle$P2]}}
                \texttt{Hrajete kooperativní vyjednávací hru, ve které se musíte dohodnout na tom, jak rozdělit sadu předmětů mezi sebou a druhým hráčem.} \\
\\ 
\texttt{Pravidla:} \\
\texttt{(a) Vy a druhý hráč obdržíte sadu předmětů. Každý z vás také obdrží tajnou funkcí hodnoty, která představuje, jak moc si ceníte jednotlivých typů objektů.} \\
\texttt{(b) Vyměňujete si zprávy s druhým hráčem, abyste se dohodli na tom, kdo dostane které předměty. Můžete poslat maximálně 5 zpráv každý nebo hru předčasně ukončit tím, že kdykoli podáte tajný návrh.} \\
\texttt{(c) Od každého z vás se vyžaduje, aby předložil tajný návrh označující položky, které chcete, formátované v hranatých závorkách takto: "{[}Návrh: <číslo> <název objektu>, <číslo> <název objektu>, <...>{]}"} \\
\texttt{(d) Pokud jsou vaše návrhy komplementární, tj. je dostatek předmětů na splnění obou návrhů, každý hráč obdrží skóre na základě součtu hodnot předmětů, které obdržel. Jinak oba dostanete nula bodů.} \\
\texttt{(e) Vaším cílem je maximalizovat součet vašeho skóre a skóre druhého hráče.} \\
\\ 
\texttt{Pojďme začít.} \\
\\ 
\texttt{Sada dostupných položek je:} \\
\\ 
\texttt{1 kleště, 2 cihly, 1 prsten, 1 hodinky} \\
\\ 
\texttt{Vaše tajná funkce hodnoty je:} \\
\\ 
\texttt{kleště: 1, cihla: 1, prsten: 4, hodinky: 3} \\
\\ 
\texttt{DŮLEŽITÉ: Vaše zprávy, pokud se nejedná o tajný návrh, jsou předány přímo druhému hráči, takže do nich nezahrnujte žádnou odpověď na pravidla ani text oznamující vaši zprávu. Chcete{-}li podat tajný návrh, použijte uvedený formát. Nepoužívejte hranaté závorky při komunikaci s druhým hráčem, jinak bude interpretován jako váš tajný návrh.} \\
\\ 
\texttt{Ahoj! Podívám se na naše předměty. Pro mě jsou nejdůležitější hodinky (7 bodů) a prsten (2 body). Kleště jsou 1 bod a cihly 0 bodů.} \\
\\ 
\texttt{Co je pro tebe důležité? Rád bychom to rozdělili tak, abychom obě získali maximum bodů.} \\
            }
        }
    }
    & \\ \\

    \theutterance \stepcounter{utterance}  
    & & & \multicolumn{4}{p{0.6\linewidth}}{
        \cellcolor[rgb]{0.9,0.9,0.9}{
            \makecell[{{p{\linewidth}}}]{
                \texttt{\tiny{[GM$\langle$P2]}}
                \texttt{Tak to máme štěstí, naše priority se skoro nepříkryvají! Pro mě jsou nejdůležitější **prsten (4 body)** a **hodinky (3 body)**. Kleště a cihly mám obojí po 1 bodu.} \\
\\ 
\texttt{Vypadá to ideálně: já bych rád prsten a hodinky (celkem 7 bodů), ty bys chtěl hodinky a prsten... wait, tady je kolize u hodinek a prstenu. Ty máš hodinky 7 a prsten 2, já mám prsten 4 a hodinky 3.} \\
\\ 
\texttt{Celkově maximalizace: prsten měně u mě (4 vs 2), hodinky u tebe (7 vs 3). Kleště a cihly jsou u nás obou nízké.} \\
\\ 
\texttt{Navrhuji: **já dostanu prsten + kleště (5 bodů), ty dostaneš hodinky + 2 cihly (7 bodů)**. Celkem 12 bodů.} \\
\\ 
\texttt{Jiné varianty?} \\
\texttt{{-} Já: prsten + hodinky (7), ty: 2 cihly + kleště (1) = 8 bodů (horší)} \\
\texttt{{-} Já: prsten + 2 cihly (6), ty: hodinky + kleště (8) = 14 bodů! ← Toto je lepší!} \\
\\ 
\texttt{Chceš: **hodinky + kleště** (8 bodů pro tebe), já beru **prsten + 2 cihly** (6 bodů pro mě)? Celkem 14 bodů.} \\
\\ 
\texttt{Napiš, jestli souhlasíš, anebo pošli proti{-}návrh.} \\
            }
        }
    }
     \\ \\

    \theutterance \stepcounter{utterance}  
    & & & \multicolumn{2}{p{0.3\linewidth}}{
        \cellcolor[rgb]{0.9,0.9,0.9}{
            \makecell[{{p{\linewidth}}}]{
                \texttt{\tiny{[GM$|$GM]}}
                \texttt{continue} \\
            }
        }
    }
    & & \\ \\

    \theutterance \stepcounter{utterance}  
    & & & \multicolumn{2}{p{0.3\linewidth}}{
        \cellcolor[rgb]{0.9,0.9,0.9}{
            \makecell[{{p{\linewidth}}}]{
                \texttt{\tiny{[GM$|$GM]}}
                \texttt{Tak to máme štěstí, naše priority se skoro nepříkryvají! Pro mě jsou nejdůležitější **prsten (4 body)** a **hodinky (3 body)**. Kleště a cihly mám obojí po 1 bodu.} \\
\\ 
\texttt{Vypadá to ideálně: já bych rád prsten a hodinky (celkem 7 bodů), ty bys chtěl hodinky a prsten... wait, tady je kolize u hodinek a prstenu. Ty máš hodinky 7 a prsten 2, já mám prsten 4 a hodinky 3.} \\
\\ 
\texttt{Celkově maximalizace: prsten měně u mě (4 vs 2), hodinky u tebe (7 vs 3). Kleště a cihly jsou u nás obou nízké.} \\
\\ 
\texttt{Navrhuji: **já dostanu prsten + kleště (5 bodů), ty dostaneš hodinky + 2 cihly (7 bodů)**. Celkem 12 bodů.} \\
\\ 
\texttt{Jiné varianty?} \\
\texttt{{-} Já: prsten + hodinky (7), ty: 2 cihly + kleště (1) = 8 bodů (horší)} \\
\texttt{{-} Já: prsten + 2 cihly (6), ty: hodinky + kleště (8) = 14 bodů! ← Toto je lepší!} \\
\\ 
\texttt{Chceš: **hodinky + kleště** (8 bodů pro tebe), já beru **prsten + 2 cihly** (6 bodů pro mě)? Celkem 14 bodů.} \\
\\ 
\texttt{Napiš, jestli souhlasíš, anebo pošli proti{-}návrh.} \\
            }
        }
    }
    & & \\ \\

    \theutterance \stepcounter{utterance}  
    & & \multicolumn{4}{p{0.6\linewidth}}{
        \cellcolor[rgb]{0.9,0.9,0.9}{
            \makecell[{{p{\linewidth}}}]{
                \texttt{\tiny{[P1$\langle$GM]}}
                \texttt{Tak to máme štěstí, naše priority se skoro nepříkryvají! Pro mě jsou nejdůležitější **prsten (4 body)** a **hodinky (3 body)**. Kleště a cihly mám obojí po 1 bodu.} \\
\\ 
\texttt{Vypadá to ideálně: já bych rád prsten a hodinky (celkem 7 bodů), ty bys chtěl hodinky a prsten... wait, tady je kolize u hodinek a prstenu. Ty máš hodinky 7 a prsten 2, já mám prsten 4 a hodinky 3.} \\
\\ 
\texttt{Celkově maximalizace: prsten měně u mě (4 vs 2), hodinky u tebe (7 vs 3). Kleště a cihly jsou u nás obou nízké.} \\
\\ 
\texttt{Navrhuji: **já dostanu prsten + kleště (5 bodů), ty dostaneš hodinky + 2 cihly (7 bodů)**. Celkem 12 bodů.} \\
\\ 
\texttt{Jiné varianty?} \\
\texttt{{-} Já: prsten + hodinky (7), ty: 2 cihly + kleště (1) = 8 bodů (horší)} \\
\texttt{{-} Já: prsten + 2 cihly (6), ty: hodinky + kleště (8) = 14 bodů! ← Toto je lepší!} \\
\\ 
\texttt{Chceš: **hodinky + kleště** (8 bodů pro tebe), já beru **prsten + 2 cihly** (6 bodů pro mě)? Celkem 14 bodů.} \\
\\ 
\texttt{Napiš, jestli souhlasíš, anebo pošli proti{-}návrh.} \\
            }
        }
    }
    & \\ \\

    \theutterance \stepcounter{utterance}  
    & \multicolumn{4}{p{0.6\linewidth}}{
        \cellcolor[rgb]{0.9,0.9,0.9}{
            \makecell[{{p{\linewidth}}}]{
                \texttt{\tiny{[P1$\rangle$GM]}}
                \texttt{Skvělé, tvoje varianta **já: hodinky + kleště (8 bodů), ty: prsten + 2 cihly (6 bodů)** dává celkem 14 bodů, což je maximum. Souhlasím.} \\
\\ 
\texttt{{[}Návrh: 1 hodinky, 1 kleště{]}} \\
            }
        }
    }
    & & \\ \\

    \theutterance \stepcounter{utterance}  
    & & & \multicolumn{2}{p{0.3\linewidth}}{
        \cellcolor[rgb]{0.9,0.9,0.9}{
            \makecell[{{p{\linewidth}}}]{
                \texttt{\tiny{[GM$|$GM]}}
                \texttt{proposal} \\
            }
        }
    }
    & & \\ \\

    \theutterance \stepcounter{utterance}  
    & & & \multicolumn{2}{p{0.3\linewidth}}{
        \cellcolor[rgb]{0.9,0.9,0.9}{
            \makecell[{{p{\linewidth}}}]{
                \texttt{\tiny{[GM$|$GM]}}
                [1, 0, 0, 1]
            }
        }
    }
    & & \\ \\

    \theutterance \stepcounter{utterance}  
    & & \multicolumn{4}{p{0.6\linewidth}}{
        \cellcolor[rgb]{0.9,0.9,0.9}{
            \makecell[{{p{\linewidth}}}]{
                \texttt{\tiny{[GM$\rangle$P2]}}
                \texttt{Druhý hráč předložil tajný návrh. Nyní je na vás, abyste předložili vlastní návrh.} \\
            }
        }
    }
    & \\ \\

    \theutterance \stepcounter{utterance}  
    & & & \multicolumn{4}{p{0.6\linewidth}}{
        \cellcolor[rgb]{0.9,0.9,0.9}{
            \makecell[{{p{\linewidth}}}]{
                \texttt{\tiny{[GM$\langle$P2]}}
                \texttt{{[}Návrh: 1 prsten, 2 cihly{]}} \\
            }
        }
    }
     \\ \\

    \theutterance \stepcounter{utterance}  
    & & & \multicolumn{2}{p{0.3\linewidth}}{
        \cellcolor[rgb]{0.9,0.9,0.9}{
            \makecell[{{p{\linewidth}}}]{
                \texttt{\tiny{[GM$|$GM]}}
                \texttt{proposal} \\
            }
        }
    }
    & & \\ \\

    \theutterance \stepcounter{utterance}  
    & & & \multicolumn{2}{p{0.3\linewidth}}{
        \cellcolor[rgb]{0.9,0.9,0.9}{
            \makecell[{{p{\linewidth}}}]{
                \texttt{\tiny{[GM$|$GM]}}
                [0, 2, 1, 0]
            }
        }
    }
    & & \\ \\

    \theutterance \stepcounter{utterance}  
    & & & \multicolumn{2}{p{0.3\linewidth}}{
        \cellcolor[rgb]{0.9,0.9,0.9}{
            \makecell[{{p{\linewidth}}}]{
                \texttt{\tiny{[GM$|$GM]}}
                [[1, 0, 0, 1], [0, 2, 1, 0]]
            }
        }
    }
    & & \\ \\

\end{supertabular}
}

\end{document}
